% Part: first-order-logic
% Chapter: model-theory
% Section: isomorphism

\documentclass[../../../include/open-logic-section]{subfiles}

\begin{document}

\olfileid[fr]{mod}{bas}{iso}
\olsection{Structures isomorphes}

Deux !!{structure}s du premier ordre peuvent se ressembler de deux
manières. Elles peuvent, d'abord, rendre vrais les mêmes !!{sentence}s.
On les dit alors \emph{élémentairement équivalentes}. Mais deux
structures peuvent être très différentes tout en rendant vrais les
mêmes !!{sentence}s : l'une peut, par exemple, être !!{enumerable}
et l'autre non. Certaines propriétés d'!!a{structure} ne sont en effet
pas exprimables dans un langage du premier ordre, soit parce que le
langage n'est pas assez riche, soit en raison de limites fondamentales
de la logique du premier ordre, comme le théorème de
L\"owenheim--Skolem. Il existe donc une manière plus exigeante de
dire que deux !!{structure}s se ressemblent : elles ne diffèrent que
par les objets qui les composent, sans différer par leurs propriétés
structurelles. La notion d'\emph{isomorphisme} précise cette idée.

\begin{defn}
\ollabel{defn:elem-equiv}
Soient deux !!{structure}s $\Struct{M}$ et $\Struct M'$ pour un même
!!{language}~$\Lang{L}$. On dit que $\Struct{M}$ et $\Struct M'$ sont
\emph{élémentairement équivalentes}, ce que l'on note
$\Struct{M} \equiv \Struct M'$, si, pour tout !!{sentence}~$!A$
de~$\Lang{L}$, $\Sat{M}{!A}$ si et seulement si $\Sat{M'}{!A}$.
\end{defn}

\begin{defn}
\ollabel{defn:isomorphism} Soient deux !!{structure}s $\Struct{M}$ et
$\Struct M'$ pour un même !!{language}~$\Lang L$. On dit que
$\Struct{M}$ et $\Struct M'$ sont \emph{isomorphes}, ce que l'on note
$\Struct{M} \simeq \Struct M'$, s'il existe une application
$h \colon \Domain{M} \to \Domain{M'}$ telle que :
\begin{enumerate}
\item $h$ est !!{injective} : si $h(x) =
  h(y)$, alors $x = y$;
\item $h$ est !!{surjective} : pour tout $y \in \Domain{M'}$, il
  existe $x \in \Domain{M}$ tel que $h(x) = y$;
\item \ollabel{defn:iso-const}pour chaque !!{constant} $c$ :
  $h(\Assign{c}{M}) = \Assign{c}{M'}$;
\item \ollabel{defn:iso-pred}pour chaque !!{predicate} $n$-aire~$P$ :
  \[
  \tuple{a_1, \dots, a_n}\in \Assign{P}{M} \quad\text{si et seulement si}\quad
  \tuple{h(a_1), \dots, h(a_n)} \in \Assign{P}{M'};
  \]
\item \ollabel{defn:iso-func}pour chaque !!{function} $n$-aire $f$ :
  \[
  h(\Assign{f}{M}(a_1, \dots, a_n)) =
  \Assign{f}{M'}(h(a_1), \dots, h(a_n)).
  \]
\end{enumerate}
\end{defn}

\begin{thm}
\ollabel{thm:isom}
Si $\Struct{M} \iso \Struct M'$, alors $\Struct{M} \elemequiv
\Struct{M'}$.
\end{thm}

\begin{proof}
Soit $h$ un isomorphisme de $\Struct{M}$ sur $\Struct M'$. Pour toute
assignation~$s$, $h \circ s$ est la composée de $h$ et de $s$,
c'est-à-dire l'assignation dans $\Struct{M'}$ telle que
$(h \circ s)(x) = h(s(x))$. Une induction sur $t$ et $!A$ permet de
prouver les assertions plus fortes suivantes :
\begin{enumerate}
  \item[a.] $h(\Value{t}{M}[s]) = \Value{t}{M'}[h\circ s]$.
  \item[b.] $\Sat{M}{!A}[s]$ si et seulement si $\Sat{M'}{!A}[h \circ s]$.
\end{enumerate}
La première se démontre par induction sur la complexité de~$t$.
\begin{enumerate}
\item Si $t \ident c$, alors $\Value{c}{M}[s] = \Assign{c}{M}$ et
  $\Value{c}{M'}[h \circ s] = \Assign{c}{M'}$. Donc
  $h(\Value{t}{M}[s]) = h(\Assign{c}{M}) = \Assign{c}{M'}$ (par
  \olref{defn:iso-const} de la \olref{defn:isomorphism}) $=
  \Value{t}{M'}[h \circ s]$.
\item Si $t \ident x$, alors $\Value{x}{M}[s] = s(x)$ et
  $\Value{x}{M'}[h \circ s] = h(s(x))$. Donc
  $h(\Value{x}{M}[s]) = h(s(x)) = \Value{x}{M'}[h \circ s]$.
\item Si $t \ident f(t_1, \dots, t_n)$, alors
  \begin{align*}
    \Value{t}{M}[s] & = \Assign{f}{M}(\Value{t_1}{M}[s], \dots, \Value{t_n}{M}[s]) \quad\text{et}\\
  \Value{t}{M'}[h \circ s] & = \Assign{f}{M'}(\Value{t_1}{M'}[h \circ
    s], \dots, \Value{t_n}{M'}[h \circ s]).
  \end{align*}
  L'hypothèse d'induction donne, pour chaque $i$,
  $h(\Value{t_i}{M}[s]) = \Value{t_i}{M'}[h\circ s]$. Ainsi,
  \begin{align}
    h(\Value{t}{M}[s])
    & = h(\Assign{f}{M}(\Value{t_1}{M}[s], \dots, \Value{t_n}{M}[s])) \notag\\
    & = \Assign{f}{M'}(h(\Value{t_1}{M}[s]), \dots,
    h(\Value{t_n}{M}[s])) \ollabel{iso-1}\\
    & = \Assign{f}{M'}(\Value{t_1}{M'}[h \circ s], \dots,
    \Value{t_n}{M'}[h \circ s]) \ollabel{iso-2}\\
    & = \Value{t}{M'}[h\circ s] \notag
  \end{align}
  L'égalité \olref{iso-1} résulte de \olref{defn:iso-func}
  dans la \olref{defn:isomorphism}, et \olref{iso-2} résulte de
  l'hypothèse d'induction.
\end{enumerate}
La démonstration de (b) est laissée en exercice.\footnote{Correction éditoriale : dans le cas du terme composé, sa valeur dans $\Struct{M'}$ est calculée avec $\Assign{f}{M'}$, et non avec $\Assign{f}{M}$ comme dans l'affichage de la source. Nous rétablissons aussi la parenthèse qui ferme l'application extérieure de $h$ dans la première ligne de la chaîne d'égalités.}

Si $!A$ est un énoncé, les assignations~$s$ et $h \circ s$ sont
sans effet sur sa valeur de vérité; on a donc $\Sat{M}{!A}$
si et seulement si $\Sat{M'}{!A}$.
\end{proof}

\begin{prob}
Démontrer en détail l'assertion (b) du
\olref[mod][bas][iso]{thm:isom}. Préciser où intervient chacune
des cinq propriétés qui caractérisent les isomorphismes dans la
\olref[mod][bas][iso]{defn:isomorphism}.
\end{prob}

\begin{defn}
Un \emph{automorphisme} d'!!a{structure}~$\Struct{M}$ est un
isomorphisme de $\Struct{M}$ sur elle-même.
\end{defn}

\begin{prob}
Montrer que, pour toute !!{structure}~$\Struct{M}$, si $X$ est une
partie définissable sans paramètres de $\Struct{M}$ et si $h$ est un automorphisme de
$\Struct{M}$, alors $X = \Setabs{h(x)}{x \in X}$ (autrement dit,
$X$ est invariant par $h$).\footnote{Précision éditoriale : nous ajoutons « sans paramètres », absent de la source. Pour une définition avec paramètres, la même conclusion est garantie si l'automorphisme fixe ces paramètres ; sans cette condition, un singleton défini avec son élément comme paramètre fournit un contre-exemple.}
\end{prob}

\end{document}
